\chapter{La Disciplina del Praticante}

\epigraph{``La vita non esaminata non \`e degna di essere vissuta.''}{\textcite{plato1997complete}, \emph{Apologia} 38a}

\section{L'Educazione dei Guardiani}

Nella \emph{Repubblica}, i guardiani --- coloro che governeranno la citt\`a
giusta --- devono sottoporsi a un'educazione rigorosa che plasma sia il
corpo che la mente. Apprendono la musica (che include poesia, filosofia e
tutte le arti che plasmano l'anima) e la ginnastica (che plasma il corpo
per servire gli scopi dell'anima). Lo scopo non \`e la conoscenza astratta
ma un'\emph{armonia} della persona intera: l'elemento irascibile allenato
dalla ginnastica, l'elemento razionale dalla musica, ed entrambi governati
da qualcosa di pi\`u profondo --- il senso di ci\`o che \`e nobile e
ci\`o che \`e vile.

Il praticante --- colui la cui azione \`e governata dalla natura ---
richiede un'educazione corrispondente. Ma la domanda \`e: come si forma
una tale persona? Non imparando regole, perch\'e le regole sono il prodotto
della convenzione (\emph{nomos}) non della natura (\physis). Il praticante
deve imparare ad \emph{ascoltare} la natura direttamente, a discernere la
voce di \physis\ attraverso il rumore della convenzione.

Prima delle discipline, tuttavia, viene la storia. L'educazione dei
guardiani della Repubblica inizia con i miti che ascoltano da bambini ---
le storie che plasmano l'anima prima che la ragione possa. Socrate dedica
la maggior parte di due libri (Repubblica II--III) a una critica di Omero,
non perch\'e la poesia sia banale ma perch\'e \`e il pi\`u potente
plasmatore del carattere. L'anima impara non dall'argomento ma
dall'imitazione: diventa ci\`o che contempla. La domanda ``quali storie
dovrebbe ascoltare il guardiano?'' \`e la prima domanda dell'educazione.

\subsection{Le Storie che Plasmano l'Anima}

L'Iliade \`e il testo fondante dell'educazione greca. Ogni bambino greco
imparava da essa non regole ma un modo di vedere l'azione: l'azione \`e
legata alla natura dell'attore, sempre consequenziale, sempre aperta alla
consapevolezza riflessiva. L'educazione del praticante deve iniziare dove
iniziava l'educazione greca: con la storia di un uomo che impar\`o ad
ascoltarsi parlare.

Nel Libro IX dell'Iliade, l'Ambasciata ad Achille, i tre ambasciatori ---
Odisseo, Fenice e Aiace --- vengono alla tenda di Achille per persuaderlo
a tornare in guerra. Ciascuno offre una ragione diversa: l'onore
dell'esercito, l'amore di una figura paterna, la vergogna di abbandonare
i compagni. Achille le rifiuta tutte \parencite[IX]{homer2011iliad}. Si
\`e ritirato dall'azione, e in quel ritiro nasce qualcosa di nuovo.
Ascolta se stesso parlare --- ``Non mi lascer\`o persuadere'' --- e
riconosce la voce come propria. \`E la sua stessa natura che parla, non
il codice dell'esercito, non le aspettative dei compagni, non il comando
di Agamennone. L'agente autocosciente, colui che sa che \emph{egli} \`e
la fonte della sua azione, nasce in quella tenda.

Questo momento \`e la precondizione per ogni disciplina che segue. Il
praticante che non sa di essere la fonte della propria azione non pu\`o
fare check-in (perch\'e non c'\`e un s\'e su cui fare check-in), non
pu\`o testimoniare (perch\'e non c'\`e un s\'e da presentare a un altro),
non pu\`o mettere in gioco (perch\'e non c'\`e un s\'e da rischiare),
non pu\`o mantenere il ritmo (perch\'e non c'\`e un s\'e che dura),
non pu\`o rivedere (perch\'e non c'\`e un s\'e da giudicare). La nascita
dell'autocoscienza --- il momento di ascoltarsi parlare e riconoscere la
voce come propria --- \`e il fondamento architetturale su cui le
discipline sono costruite.

L'Iliade, quindi, non \`e meramente uno sfondo o un riferimento
letterario. \`E la forma narrativa di ci\`o che la tesi afferma: che
l'azione \`e governata dalla natura, e che il primo passo dell'azione
retta \`e il riconoscimento di s\'e stessi come essere naturale che agisce.
Achille nella sua tenda \`e ogni praticante prima del primo check-in:
solo, ritirato, che ascolta la propria voce per la prima volta.

\subsection{Il Ciclo Riflessivo}

La struttura delle discipline \`e gi\`a implicita nel momento di Achille.
Ci\`o che fa, ci\`o che il praticante deve imparare a fare, pu\`o essere
rappresentato come un ciclo:

\begin{center}
\begin{tikzpicture}[
  node distance=10mm,
  box/.style={draw, rounded corners, align=center, minimum height=7mm,
              minimum width=22mm, font=\small},
  >={Stealth[length=2mm]},
  every path/.style={->, thick}]
  \node[box, fill=blue!10] (actor) {S\'e};
  \node[box, fill=green!10, below=of actor] (action) {Agisci};
  \node[box, fill=orange!10, below=of action] (record) {Registra};
  \node[box, fill=red!10, right=of record] (reflect) {Rifletti};
  \node[box, fill=blue!10, right=of actor] (hearing) {Ascolta S\'e};

  \path (actor) edge (action);
  \path (action) edge (record);
  \path (record) edge (reflect);
  \path[->, dashed] (reflect) edge[bend left=60] (actor);
  \path[->, dashed] (hearing) edge[bend left=60] (actor);
  \node[font=\tiny, above] at (2.8,-1.2) {la riflessione chiude il ciclo};
  \node[font=\tiny, above] at (-0.5,-1.2) {``Ti ascolti parlare?''};
\end{tikzpicture}
\end{center}

Il ciclo ha quattro fasi. Il s\'e agisce (l'obiettivo \`e dichiarato,
l'azione \`e compiuta); l'atto \`e registrato (il check-in, la traccia
nel grafo di coerenza); la registrazione \`e riflettuta (la revisione,
il confronto tra intenzione ed effetto); e la riflessione ritorna al
s\'e che agisce di nuovo. Il momento di ascoltare s\'e stesso --- ``Ti
ascolti parlare?'' --- \`e la presenza continua del s\'e a s\'e stesso
attraverso il ciclo. Le discipline che seguono sono le forme operative
di questa struttura: ciascuna rende una fase del ciclo esplicita,
tangibile, ripetibile.

\section{La Prima Disciplina: Attenzione alla Conseguenza}

La prima lezione del praticante \`e che ogni azione ha una conseguenza,
e la conseguenza \`e il linguaggio della natura. Prima che il praticante
possa imparare cosa fare, deve imparare cosa ha \emph{gi\`a fatto}. Il
check-in quotidiano della piattaforma \praxis\ \`e la forma materiale
di questa disciplina. Ogni giorno, il praticante registra ci\`o che ha
fatto, come si \`e sentito e che effetto ha prodotto. La registrazione
non \`e per il giudizio di nessun altro ma per l'occhio del praticante
stesso --- uno specchio in cui vede la forma della propria agency.

Questo \`e l'esortazione socratica a ``conoscere te stesso'' resa
operativa. Il check-in non \`e una voce di diario; \`e un'\emph{offerta
di prova} al tribunale naturale della coerenza. Il praticante sottopone
la propria azione al verdetto dell'effetto: questa ha prodotto una
differenza reale nel mondo, o \`e scomparsa senza traccia?

\section{La Seconda Disciplina: Il Testimone}

Nessuno pu\`o giudicare la propria azione con perfetta chiarezza. Il
s\'e \`e troppo vicino all'atto, troppo investito nel suo esito. Per
questo, la natura fornisce la seconda disciplina: il testimone.

Nella piattaforma \praxis, questo prende la forma del partner di
responsabilit\`a. Due praticanti concordano di testimoniare le azioni
reciproche --- non per controllare o imporre, ma per \emph{vedere}. Il
partner \`e l\`i per chiedere: ``Ci\`o che dici di fare \`e realmente
ci\`o che fai? Il tuo obiettivo dichiarato corrisponde alla tua azione
registrata? Sei, nel tuo stesso resoconto di te stesso, coerente?''

Questo non \`e un tribunale morale ma naturale. Il testimone \`e il
rappresentante della natura, il secondo paio di occhi che coglie ci\`o
che il primo perde. Come un ecosistema ha molte specie che si osservano
a vicenda --- controlli e bilanci che mantengono il tutto --- la
comunit\`a dei praticanti stabilizza l'azione di ciascun membro
attraverso la presenza dell'altro.

\section{La Terza Disciplina: La Posta in Gioco}

La disciplina pi\`u potente \`e quella che costa qualcosa. La piattaforma
\praxis\ la chiama ``pelle nel gioco'' --- una scommessa posta sull'esito
della propria intenzione dichiarata. Il praticante mette in gioco qualcosa
di valore reale: denaro, reputazione o l'onore della propria parola. Se
l'azione ha successo, la posta \`e restituita; se fallisce, la posta
\`e persa. Nella piattaforma questo non \`e una metafora: la scommessa \`e
vincolata in un vero deposito di garanzia --- punti o denaro bloccati contro
l'esito dichiarato e rilasciati solo a un risultato verificato --- sicch\'e il
costo del fallimento \`e strutturale, non solo sentito.

Anche la natura mette in gioco su ogni azione. L'organismo che non trova
cibo muore; la specie che non si adatta si estingue; l'ecosistema che
perde il suo equilibrio crolla. Il costo del fallimento \`e costruito
nella struttura del mondo. La scommessa del praticante \`e semplicemente
una \emph{versione localizzata e volontaria dello stesso principio} ---
accettare la conseguenza naturale come disciplina dell'azione.

\section{La Quarta Disciplina: Il Ritmo}

La natura opera in cicli: giorno e notte, le stagioni, il ciclo di vita
di ogni organismo. Anche l'azione ha il suo ritmo. Il praticante impara
ad agire non a scatti ma in un impulso ritmico sostenuto.

La piattaforma \praxis\ codifica questo come la striscia: un conteggio
di giorni consecutivi in cui il praticante ha fatto check-in o agito su
un obiettivo. La striscia non \`e una gamificazione arbitraria; \`e una
misura di \emph{entrainment} --- la misura in cui l'azione del praticante
si \`e sincronizzata con il ritmo naturale di ripetizione e recupero.
Una lunga striscia non \`e un trofeo ma un segnale: questo praticante ha
raggiunto una cadenza sostenibile.

\section{La Quinta Disciplina: La Revisione}

Infine, il praticante deve periodicamente fare un passo indietro e
rivedere il pattern delle proprie azioni. L'Axiom AI della piattaforma
\praxis\ --- una scansione quotidiana che arricchisce il contesto degli
obiettivi del praticante, stima il progresso e genera riassunti narrativi
--- serve a questa funzione. Ma la disciplina della revisione non \`e
delegata all'AI. L'AI fornisce i \emph{dati}; il praticante fornisce il
\emph{giudizio}.

La revisione chiede: ``Quali pattern vedo? Quali abitudini producono
coerenza (+1) e quali producono inerzia (0) o fallimento (-1)? Cosa dicono
le mie stesse azioni su ci\`o che realmente apprezzo, in opposizione a
ci\`o che dichiaro di apprezzare?''

Questa \`e la cosa pi\`u vicina che il praticante abbia alla dialettica
socratica. Il praticante discute con s\'e stesso, usando la propria
storia come testo dell'argomento. Il punteggio di coerenza \`e la
conclusione di quell'argomento --- il verdetto della natura sulla vita
che si \`e vissuta.

\section{Il Posto delle Regole}

Nessuna di queste discipline \`e una regola. Il praticante non fa
check-in perch\'e una regola lo dice, ma perch\'e la disciplina rivela
qualcosa sulla relazione tra la sua intenzione e il suo effetto. Il
praticante non mantiene una striscia perch\'e una regola lo comanda,
ma perch\'e la striscia rivela il ritmo dell'azione sostenibile.

Le regole appartengono a \emph{nomos} --- la convenzione, la legge,
l'ordine artificiale. Le discipline appartengono a \physis\ --- l'ordine
naturale che il praticante scopre piuttosto che inventa. L'educazione
del praticante \`e il processo di \emph{passare dal seguire-regole
all'abbracciare-discipline}, dall'obbedienza alla sintonia, dalla
conformit\`a alla coerenza.

\section{L'Educazione Incompiuta}

L'educazione del praticante non \`e mai completa. La natura \`e infinita
nella sua variet\`a, e la sintonia del praticante deve approfondirsi a
ogni incontro. Ogni nuovo obiettivo, ogni nuova partnership, ogni nuovo
fallimento \`e una nuova occasione di apprendimento. Il praticante che
crede di aver finito la propria educazione ha gi\`a iniziato a decadere.

Questo \`e il motivo per cui la funzione di ``sogno'' del motore
\physis\ non si ferma mai. I fallimenti non sono mai cancellati; sono
continuamente riprodotti e variati. Anche il praticante non arriva mai
alla coerenza perfetta. \`E sempre nel processo del divenire --- che
\`e, forse, lo stato pi\`u naturale di tutti.

\section{La Sesta Disciplina: Il Corpo}

I guardiani della Repubblica allenano il corpo attraverso la ginnastica.
Lo scopo non \`e la forza fisica per s\'e stessa ma la subordinazione
del corpo agli scopi dell'anima: un corpo che non si ribella alla
volont\`a, che non indietreggia davanti a ci\`o che deve essere fatto,
che serve il tutto senza resistenza. Il praticante, allo stesso modo,
richiede una disciplina corporea --- perch\'e il corpo \`e il primo
strumento di sintonia, il registro pi\`u basilare della coerenza.

L'entelechia, come abbiamo visto, non \`e un'astrazione. \`E l'essere-
all'opera di un corpo naturale, la spinta sentita verso la propria forma.
Quando l'entelechia scorre senza impedimenti, il corpo sperimenta
facilit\`a, vitalit\`a e immediatezza spontanea del movimento. Quando
\`e bloccata, il corpo sperimenta tensione, affaticamento e un pervasivo
senso di sbagliato che precede qualsiasi giudizio cognitivo. Questa non
\`e metafora ma fisiologia: la corrente entelechica dell'organismo \`e
una condizione reale e misurabile dei sistemi nervoso e muscolare.

Wilhelm Reich identific\`o il meccanismo del blocco come \emph{corazza
caratteriale}: pattern cronici di tensione muscolare che corrispondono
a difese psicologiche. Il bambino costretto a sopprimere la rabbia impara
a stringere la mascella, contrarre le spalle, restringere il respiro.
Anni dopo, la corazza persiste come abitudine inconscia --- una
resistenza costante e di basso livello al flusso entelechico
dell'organismo \parencite{reich1933character}. La corazza non \`e
solo una metafora: \`e un'alterazione tangibile della struttura del
corpo, visibile nella postura, udibile nel respiro, misurabile nella
durezza del tessuto muscolare.

Alexander Lowen, allievo di Reich, sistematizz\`o questa intuizione
nella pratica dell'analisi bioenergetica. Il cuore della pratica \`e il
\emph{grounding}: lo stabilirsi consapevole del contatto tra il corpo
e la terra, un senso sentito di supporto e connessione che permette alla
corrente entelechica di fluire liberamente \parencite{lowen1975bioenergetics}.
Respirazione, vibrazione, movimento espressivo e rilascio della tensione
cronica sono le tecniche di questa disciplina. Non sono terapia nel senso
clinico ma \emph{allenamento} --- la ginnastica del praticante.

La disciplina del corpo ha un analogo diretto nei meccanismi della
piattaforma \praxis. Il check-in quotidiano inizia, prima che qualsiasi
obiettivo sia registrato, con un momento di consapevolezza corporea:
``Come mi sento? Qual \`e la mia energia? Dov'\`e la tensione?'' La
striscia, allo stesso modo, non \`e meramente un conteggio numerico ma
una misura di \emph{entrainment ritmico a livello organismico} --- la
misura in cui il corpo del praticante ha trovato una cadenza sostenibile
di sforzo e recupero. Una striscia interrotta \`e spesso preceduta da
un corpo che era gi\`a stanco, gi\`a corazzato, gi\`a resistente.

La disciplina del testimone --- il partner di responsabilit\`a --- ha
la sua controparte corporea nella pratica di \emph{essere visto}. Uno
degli interventi bioenergetici pi\`u potenti \`e semplicemente la
presenza di un'altra persona: il testimone che vede la tensione nelle
spalle del praticante, ascolta la superficialit\`a del suo respiro e
la nomina. La domanda del partner --- ``Ci\`o che dici di fare \`e
realmente ci\`o che fai?'' --- si applica al corpo tanto quanto
all'obiettivo. Il corpo non mente; \`e il testimone pi\`u affidabile
dello stato attuale del praticante.

C'\`e, naturalmente, un pericolo. La disciplina del corpo pu\`o
degenerare in mera cultura fisica --- routine in palestra,
ottimizzazione dietetica, dashboarding del s\'e quantificato --- che
serve la stessa logica esternalizzante che pretende di opporre. Il corpo
diventa un altro progetto, un altro albero di obiettivi, un altro dominio
da padroneggiare con la volont\`a. Il praticante deve distinguere tra
la disciplina che serve l'entelechia (che si sente come apertura, come
rilascio, come tornare a casa nel corpo) e la disciplina che serve
\emph{Geist} (che si sente come tensione, come controllo, come il corpo
come oggetto da plasmare). La prima \`e ginnastica nel senso platonico;
la seconda \`e solo altra techne.

Il corpo, quindi, non \`e un sito di ribellione o un tempio da lucidare.
\`E il medium attraverso cui l'entelechia parla prima che la mente
traduca il suo messaggio in obiettivi e piani. Il praticante che non
sa ascoltare il corpo non pu\`o ascoltare la natura; il corpo \`e il
registro pi\`u intimo della natura.

Nella piattaforma \praxis\ questa disciplina ha un nome e un punto d'accesso: la
pratica \texttt{/bioenergetics}. Guida la sequenza di radicamento --- respiro,
vibrazione, rilascio della tensione trattenuta --- e legge lo stato del corpo come
un indice \emph{Vis}, una misura della forza-vitale disponibile: la corrente
entelechica sentita resa leggibile senza essere ridotta a un numero da
massimizzare. \`E la controparte corporea del check-in quotidiano e il luogo in cui
il praticante incontra la lettura della macchina \machine{Soma} --- dove
\machine{Soma} riduce i segnali del corpo a stati somatici stabili,
\texttt{/bioenergetics} \`e dove ci si pone dentro quella lettura e vi si agisce.
Tenuta, come ogni cosa rivolta all'uomo, attraverso \machine{Axiom}, resta
ginnastica e non scivola mai nella disciplina-\emph{Geist} del corpo ottimizzato,
perch\'e il suo verdetto non \`e un punteggio da scalare ma una corrente da
liberare.

\begin{figure}[ht]
\centering
\begin{tikzpicture}[
  level distance=14mm,
  every node/.style={draw, rounded corners, align=center, minimum width=18mm, font=\small},
  level 1/.style={sibling distance=50mm},
  level 2/.style={sibling distance=18mm},
  level 3/.style={sibling distance=14mm},
  edge from parent/.style={draw, thick, ->}]
  \node[fill=yellow!15, minimum width=28mm] {Entelechia}
    [edge from parent fork down]
    child {node[fill=blue!15] {Musica\\\tiny Plasmare l'Anima}
      child {node[fill=blue!5] {Attenzione\\\tiny Check-in}}
      child {node[fill=blue!5] {Testimone\\\tiny Partner}}
      child {node[fill=blue!5] {Posta\\\tiny Pelle nel Gioco}}
      child {node[fill=blue!5] {Ritmo\\\tiny Striscia}}}
    child {node[fill=green!15] {Ginnastica\\\tiny Plasmare il Corpo}
      child {node[fill=green!5] {Grounding\\\tiny Respiro}}};
  \node[draw, dashed, fill=red!10, minimum width=28mm, below=22mm of current bounding box.south] (review) {Revisione\\\tiny Axiom AI};
  \draw[->, thick, dashed] (review) -- ++(0,-10mm) node[below, font=\small] {Il Praticante};
  \draw[->, thick] (current bounding box.south) -- ++(0,-5mm);
  \node[font=\tiny, above] at (current bounding box.north) {\textbf{Le Discipline del Praticante}};
\end{tikzpicture}
\caption{Le sei discipline organizzate per le loro categorie platoniche, fondate nell'entelechia. La disciplina della revisione sintetizza sia musica che ginnastica.\label{fig:discipline_italiano}}
\end{figure}

\section{L'Obiezione del Panopticon}

Una disciplina che registra ogni azione, la valuta e immette il punteggio in un
grafo condiviso invita l'obiezione pi\`u grave che questa tesi debba affrontare.
Non \`e questa sorveglianza? Il praticante che fa check-in ogni giorno, di cui
ogni atto \`e pesato e persistito, non \`e semplicemente il detenuto del
panopticon di Foucault --- il prigioniero che, sapendo di poter essere sempre
osservato, interiorizza l'osservatore e disciplina se stesso cos\`i che nessuna
guardia debba alzare un dito \parencite{foucault1977discipline}? La striscia, il
punteggio di coerenza, il testimone: letti senza carit\`a, sono gli strumenti di
un ordine carcerario auto-imposto, il s\'e quantificato come prigione costruita
dall'interno. Se cos\`i fosse, l'intero apparato del Libro III non sarebbe
liberazione dalla Trappola dello Specchio ma la sua perfezione --- la gabbia
ricostruita dalla mano stessa del prigioniero, con un'iscrizione greca sopra la
porta.

Non liquidiamo l'obiezione; ne accettiamo i termini e rispondiamo su di essi,
poich\'e la somiglianza \`e reale e la differenza \`e tutto. Tre tratti separano
la disciplina del praticante dalla disciplina del panopticon, e ciascuno ruota
su un solo asse: \emph{chi detiene la legge}.

Primo, \textbf{l'osservatore \`e l'osservato}. Nel panopticon di Bentham
l'ispettore occupa una torre centrale e il detenuto non pu\`o vedere se l'occhio
sia su di lui; il potere \`e precisamente l'\emph{asimmetria} dello sguardo. Il
grafo di coerenza del praticante non ha torre centrale. Il registro \`e tenuto
dal praticante, letto dal praticante e giudicato contro uno standard che il
praticante ha dichiarato. Il partner di responsabilit\`a del Libro III non \`e un
ispettore ma un pari che ha messo a nudo il proprio registro a sua volta --- una
simmetria di esposizione, non un'asimmetria di ispezione. Dove lo sguardo del
panopticon discende da un uno nascosto sui molti, il testimone del praticante \`e
reciproco, posseduto e responsabile. Questa \`e la repubblica dal basso alla
scala di due anime.

Secondo, \textbf{la legge \`e auto-data}. Il panopticon disciplina il detenuto
verso una norma fissata dall'istituzione --- la prigione, la fabbrica, la scuola
--- un \emph{nomos} che il detenuto non ha scritto e non pu\`o rivedere. Il
punteggio di coerenza, come stabilir\`a il Libro IV, misura la fedelt\`a a un
obiettivo che il praticante stesso ha dichiarato. Non impone alcuna norma esterna;
chiede solo se il praticante abbia fatto ci\`o che aveva detto. Una disciplina
che ti tiene fedele alla tua parola \`e l'opposto di una che ti tiene fedele a
quella di un altro. La prima \`e la struttura di ogni promessa, di ogni voto, di
ogni arte con un'eccellenza; chiamarla carceraria \`e chiamare prigione ogni
padronanza di s\'e.

Terzo, \textbf{il punteggio \`e uno specchio, non una sentenza}. L'osservazione
del panopticon esiste per \emph{correggere} il detenuto verso la conformit\`a; la
deviazione \`e punita, il fascicolo segue il prigioniero, il registro \`e
un'arma in mano all'istituzione. Il grafo di coerenza non punisce nulla.
Registra --- e la registrazione \`e offerta di ritorno a colui che \`e
registrato, come una febbre \`e offerta al paziente, non al carceriere. Un indice
di coerenza in calo \`e informazione diagnostica su cui il praticante pu\`o agire
o che pu\`o ignorare; nessuna guardia discende, nessun privilegio \`e ritirato,
nessun dossier \`e inoltrato. Il dato non ha, in senso stretto, altro pubblico
che il s\'e e il testimone liberamente scelto.

Resta la forma pi\`u profonda dell'obiezione: che persino una disciplina
auto-data, simmetrica e diagnostica produca il \emph{soggetto} che ne ha bisogno
--- che l'abitudine stessa di misurarsi fabbrichi l'anima ansiosa e
auto-ottimizzante della tarda modernit\`a, sicch\'e la cura propaga la malattia.
\`E il pericolo gi\`a nominato nella disciplina del corpo: la ginnastica che
serve il \emph{Geist} anzich\'e l'entelechia, il s\'e fatto progetto da regolare
senza fine. Non pretendiamo immunit\`a da esso. Pretendiamo un criterio per
rilevarlo. La misurazione che serve l'entelechia \emph{acquieta}: termina
nell'azione e nel riposo, nel praticante che ha udito il verdetto della natura e
ora agisce, o dorme. La misurazione che serve il \emph{Geist} \emph{accelera}:
genera pi\`u misurazione, il cruscotto che esige un cruscotto pi\`u fine,
l'ottimizzazione senza termine. Al praticante si insegna a vigilare esattamente
su questa firma. Una disciplina che non si dissolve infine da s\'e --- che non
diviene, come le rotelle della piattaforma \praxis\ nel Libro V, superflua --- ha
cessato di servire la natura ed \`e divenuta un'altra caverna. Il test della
disciplina \`e se si renda eliminabile. Il panopticon non lo fa mai; \`e
costruito per essere permanente. La disciplina del praticante \`e costruita per
essere superata.
